\documentclass[12pt]{article}

\usepackage[spanish]{babel}
\usepackage[margin=2.5cm]{geometry}
\usepackage{setspace}

\begin{document}

\begin{center}
    {\LARGE \textbf{SOLICITUD DE PROPUESTA (RFP)}}\\[0.5cm]
    {\large \textbf{Implementación de Hub 360º con IA para la Gestión Integral de Contratos y Cumplimiento de Proveedores}}
\end{center}

\vspace{1cm}

\begin{tabular}{ll}
\textbf{Referencia de Licitación:} & RFP-GB-TI-2026-0045 \\
\textbf{Empresa Emisora:} & El Cuino S.A. de C.V. \\
\textbf{Fecha de Emisión:} & 24 de junio de 2026 \\
\end{tabular}

\vspace{1cm}

\noindent\rule{\textwidth}{0.4pt}

\vspace{0.5cm}

% ==================== 1. Introduction ====================

\section{INTRODUCCIÓN Y CONFIDENCIALIDAD}
\subsection{Propósito del Documento}
El presente documento tiene como objetivo invitar a empresas especializadas en
Inteligencia Artificial y automatización de procesos a presentar una propuesta técnica y
comercial para el desarrollo e implementación de un "Hub 360º de Gestión de
Proveedores". Este documento establece los lineamientos, alcances, requerimientos
técnicos, legales y criterios de evaluación bajo los cuales El Cuino tomará la decisión
de adjudicación.

\subsection{Acuerdo de Confidencialidad}
Toda la información contenida en este RFP, así como los datos operativos, volúmenes
de compra y procesos internos compartidos durante la fase de preguntas, es
estrictamente confidencial. La participación en esta licitación constituye la aceptación
tácita del Acuerdo de Confidencialidad previamente firmado por las partes.

% ==================== 2. Reglas y Calendario ====================

\section{REGLAS DE LA LICITACIÓN Y CALENDARIO}
El proceso de licitación se regirá por el siguiente calendario estricto. Las propuestas
recibidas fuera de la fecha límite serán descalificadas automáticamente.
\begin{itemize}
    \item Emisión del RFP: 24 de Junio de 2026.
    \item Recepción de dudas y preguntas: Hasta el 01 de Julio de 2026.
    \item Respuestas a dudas por parte de Bafar: 03 de Julio de 2026.
    \item Fecha Límite para Entrega de Propuestas: 06 de Julio de 2026 (18:00 hrs, CDMX).
    \item Presentación y Defensa (Pitch) de Finalistas: 03 al 07 de Agosto de 2026.
    \item Fallo y Adjudicación del Proyecto: 14 de Agosto de 2026.
\end{itemize}

% ==================== 3. Contexto y Problematica ====================

\section{CONTEXTO DEL NEGOCIO Y PROBLEMÁTICA
ACTUAL}
\subsection{Visión General}
El Cuino administra actualmente una red de más de 5,000 proveedores activos (materia
prima, logística, tecnología, servicios y mantenimiento industrial). Dado nuestro giro en
la industria alimentaria, operamos bajo marcos regulatorios críticos (TIF, FDA,
SENASICA, ISO 22000, FSSC 22000).

\subsection{Puntos de Dolor (Pain Points)}
\begin{enumerate}
    \item \textbf{Auditoría Manual:} El equipo de compras y legal invierte más de 4,000 horas anuales en leer, clasificar y extraer información de contratos, anexos y SLAs (Acuerdos de Nivel de Servicio).
    \item \textbf{Riesgo de Cumplimiento:} Se han detectado casos donde proveedores han operado con certificaciones de inocuidad vencidas debido a la falta de trazabilidad, exponiendo al Grupo a multas o cierres de líneas de producción.
    \item \textbf{Fugas Financieras:} La dispersión de contratos impide consolidar los días de crédito y descuentos por pronto pago. El sistema actual (ERP) no bloquea pagos si el proveedor está en incumplimiento contractual.
\end{enumerate}

% ==================== 4. Alcance del Proyecto ====================

\section{ALCANCE DEL PROYECTO (SCOPE OF WORK)}
El proveedor adjudicado deberá entregar una solución integral ("Hub 360º") basada en
algoritmos de Procesamiento de Lenguaje Natural (NLP) y Visión Computacional (OCR
avanzado). El alcance incluye:
\begin{itemize}
\item Diseño, entrenamiento y despliegue del modelo de IA.
\item Integración mediante APIs con el ecosistema actual de Bafar (SAP ERP y
Google Workspace).
\item Capacitación a 150 usuarios clave (Key Users).
\item Soporte técnico y mantenimiento evolutivo por 24 meses post-salida a
producción (Go-Live).
\end{itemize}

% ==================== 5. Requerimientos Funcionales ====================

\section{REQUERIMIENTOS FUNCIONALES OBLIGATORIOS}
\textbf{Módulo 1: Ingesta y Extracción Inteligente (NLP/OCR)}
\begin{itemize}
    \item Capacidad de ingestar contratos en formato PDF (escaneados y nativos), Word
    e imágenes.
    \item Extracción automática de: Entidades legales, fechas de inicio y fin, montos,
    penalizaciones, días de crédito y cláusulas de terminación anticipada.
    \item El modelo debe tener un nivel de confianza (Accuracy) mínimo del 92% en la
extracción de datos críticos desde el primer mes.
\end{itemize}
\textbf{Módulo 2: Auditoría de Certificaciones y Riesgo Operativo}
\begin{itemize}
    \item Motor de IA capaz de leer certificados emitidos por terceros (TIF, SENASICA) y
    contrastar las fechas de validez contra la base de datos de proveedores activos.
    \item Generación de un "Semáforo de Riesgo" (Verde, Amarillo, Rojo) por proveedor
    actualizado en tiempo real.
\end{itemize}
\textbf{Módulo 3: Flujos de Acción y "Muerte Súbita"}
\begin{itemize}
    \item Regla de Negocio Crítica: Si la IA detecta que la certificación sanitaria
    obligatoria de un proveedor ha expirado, debe enviar una instrucción automática
al ERP (vía API) para bloquear la emisión de cualquier pago o factura hasta
que el documento actualizado sea ingresado y validado por el algoritmo.
\end{itemize}
\textbf{Módulo 4: Notificaciones e Integración}
\begin{itemize}
    \item Sistema de alertas escalonadas (90, 60, 30 días antes del vencimiento) enviadas
    al correo electrónico del gestor interno y del proveedor.
\end{itemize}

% ==================== 6. Requerimientos Tecnicos ====================

\section{REQUERIMIENTOS TÉCNICOS Y DE SEGURIDAD
DE LA INFORMACIÓN}
El tratamiento de datos corporativos es de máxima sensibilidad. La arquitectura
tecnológica de la propuesta debe cumplir estrictamente con:
\begin{itemize}
    \item \textbf{Hospedaje y Cloud:} La solución debe estar alojada en una nube privada virtual
(VPC) dedicada exclusivamente a El Cuino (ej. AWS, Azure o GCP), con
residencia de datos en la región Norteamérica.
    \item \textbf{Cifrado:} Toda la información en tránsito debe estar cifrada bajo protocolo TLS
1.2 o superior. La información en reposo (contratos) debe usar cifrado AES-256.
    \item \textbf{Autenticación:} Integración con el Directorio Activo de El Cuino vía SAML 2.0 /
OAuth2 para Single Sign-On (SSO).
    \item \textbf{Cumplimiento:} El proveedor debe presentar evidencia de certificación ISO
27001 o dictamen SOC 2 Tipo II vigente.    
\end{itemize}

% ==================== 7. Estructura ====================

\section{ESTRUCTURA REQUERIDA DE LA PROPUESTA}
Los postores deberán estructurar su respuesta en el siguiente orden. La omisión de
alguna sección será motivo de descalificación:
\begin{enumerate}
    \item \textbf{Resumen Ejecutivo: (Máx. 2 páginas)}.
    \item \textbf{Perfil de la Empresa Proveedora:} Experiencia, clientes de referencia en el
    sector retail/manufactura, estructura organizacional.
    \item \textbf{Propuesta Técnica y Arquitectura IA:} Explicación detallada del modelo de IA,
    frameworks a utilizar (ej. Transformers, LLMs), arquitectura de integración y
    flujos de datos.
    \item \textbf{Diseño Experimental (Prueba de Concepto - PoC):} Descripción de cómo se
    validará el porcentaje de Accuracy y la reducción de falsos positivos en la
    extracción de contratos. (Nota: Aquí es donde tu equipo encaja el entregable
    técnico de la maestría).
    \item \textbf{Plan de Proyecto y Cronograma (Gantt):} Desglose de fases, hitos y tiempos
    estimados de implementación.
    \item \textbf{Equipo de Trabajo:} CVs del líder de proyecto, arquitectos de IA y consultores
    asignados.
    \item \textbf{Propuesta Económica (Oculta en anexo separado):} Desglose de costos de
    licencias (si aplica), servicios profesionales, infraestructura cloud y soporte. TCO
(Total Cost of Ownership) a 3 años.
\end{enumerate}

% ==================== 8. Criterios de Evaluacion ====================

\section{CRITERIOS DE EVALUACIÓN}
El Comité de Adquisiciones de TI evaluará las propuestas bajo una matriz ponderada
de 100 puntos:
\begin{itemize}
    \item \textbf{Adecuación Técnica de la Solución IA:} 40% (Precisión del modelo, facilidad
    de integración y usabilidad).
    \item \textbf{Propuesta Económica (TCO):} 25%
    \item \textbf{Experiencia y Casos de Éxito del Proveedor:} 20%
    \item \textbf{Metodología de Implementación y Tiempos:} 10%
    \item \textbf{Nivel de Soporte y Seguridad de la Información:} 5%
\end{itemize}

\end{document}